\documentclass[10pt,a4paper]{article}

\usepackage[a4paper,top=13mm,bottom=14mm,left=14mm,right=14mm]{geometry}
\usepackage{fontspec}
\usepackage{xcolor}
\usepackage{tabularx}
\usepackage{array}
\usepackage[most]{tcolorbox}
\usepackage{fancyhdr}
\usepackage{enumitem}
\usepackage{amssymb}
\usepackage{microtype}

\IfFontExistsTF{Arial}{
  \setmainfont{Arial}
  \setsansfont{Arial}
}{
  \setmainfont{TeX Gyre Heros}
  \setsansfont{TeX Gyre Heros}
}

\definecolor{AIEPRed}{HTML}{D62027}
\definecolor{AIEPNavy}{HTML}{082B5C}
\definecolor{AIEPInk}{HTML}{182B3A}
\definecolor{AIEPMuted}{HTML}{5F6B7A}
\definecolor{AIEPLine}{HTML}{C9D2DE}
\definecolor{AIEPSoft}{HTML}{F5F2EC}
\definecolor{AIEPBlueSoft}{HTML}{E9F1F8}
\definecolor{AIEPCyan}{HTML}{35B7C6}
\definecolor{AIEPGold}{HTML}{D8A928}
\definecolor{AIEPGreen}{HTML}{2E8B57}

\pagestyle{fancy}
\fancyhf{}
\renewcommand{\headrulewidth}{0pt}
\fancyfoot[L]{\footnotesize\textcolor{AIEPMuted}{GeoGreen Escolar · Taller 3}}
\fancyfoot[R]{\footnotesize\textcolor{AIEPMuted}{Ficha de propuesta tecnológica · \thepage/4}}

\setlength{\parindent}{0pt}
\setlength{\parskip}{3pt}
\renewcommand{\arraystretch}{1.16}
\setlist[itemize]{leftmargin=*,topsep=1pt,itemsep=1pt}
\newcolumntype{Y}{>{\raggedright\arraybackslash}X}
\newcolumntype{C}[1]{>{\centering\arraybackslash}p{#1}}

\newtcolorbox{fieldbox}[2][]{
  colback=white,colframe=AIEPLine,boxrule=0.55pt,arc=1mm,
  left=2.5mm,right=2.5mm,top=2mm,bottom=2mm,before skip=2pt,after skip=2pt,
  title={\scriptsize\bfseries\MakeUppercase{#2}},coltitle=AIEPNavy,fonttitle=\sffamily,
  attach title to upper={\par\vspace{0.8mm}},#1
}
\newtcolorbox{accentfield}[2][]{
  colback=AIEPSoft,colframe=AIEPRed,boxrule=0.7pt,arc=1mm,
  left=2.5mm,right=2.5mm,top=2mm,bottom=2mm,before skip=2pt,after skip=2pt,
  title={\scriptsize\bfseries\MakeUppercase{#2}},coltitle=AIEPRed,fonttitle=\sffamily,
  attach title to upper={\par\vspace{0.8mm}},#1
}
\newcommand{\writearea}[1]{\vspace{#1}\hrule height 0.35pt\vspace{4mm}\hrule height 0.35pt}
\newcommand{\checkline}[1]{$\square$\hspace{1.5mm}#1}
\newcommand{\pagehead}[2]{%
  {\small\bfseries\textcolor{AIEPRed}{GEOGREEN ESCOLAR · TALLER 3}}\par
  \vspace{2mm}
  {\LARGE\bfseries\color{AIEPNavy}#1\par}
  {\normalsize\color{AIEPMuted}#2\par}
  \vspace{4mm}
}

\begin{document}

% 1 ---------------------------------------------------------------------------
\pagehead{De una necesidad a una idea}{Ficha de propuesta tecnológica inicial · Completar en equipo}

\begin{tabularx}{\textwidth}{@{}>{\bfseries\color{AIEPNavy}}p{23mm}Y >{\bfseries\color{AIEPNavy}}p{29mm}p{40mm}@{}}
Equipo & \rule{58mm}{0.35pt} & Proyecto & \rule{38mm}{0.35pt}\\
Curso/bloque & \rule{58mm}{0.35pt} & Fecha & \rule{38mm}{0.35pt}\\
\end{tabularx}

\begin{accentfield}{La idea en una frase}
Queremos crear \rule{55mm}{0.35pt} para ayudar a \rule{42mm}{0.35pt}
porque \rule{55mm}{0.35pt}.
\vspace{4mm}
\end{accentfield}

\begin{fieldbox}[height=43mm]{Problema y contexto}
¿Qué situación quieren abordar? ¿Dónde y cuándo ocurre? Descríbanla con hechos observables, no solo con una opinión.
\writearea{14mm}
\end{fieldbox}

\begin{fieldbox}[height=35mm]{Personas o comunidad afectada}
¿Quién vive esta situación? ¿Qué cambia para esas personas si el problema mejora?
\writearea{9mm}
\end{fieldbox}

\begin{fieldbox}[height=38mm]{Punto de partida}
¿Qué recuperan del Taller 1 y del Taller 2? Anoten el problema priorizado y el dato más útil sobre el material o residuo analizado.
\writearea{11mm}
\end{fieldbox}

\begin{fieldbox}[height=35mm]{Qué hará propia nuestra propuesta}
Puede ser el contexto, la forma de medir, la respuesta, la experiencia de uso o la combinación de componentes. No basta con copiar GeoGreen.
\writearea{9mm}
\end{fieldbox}

\begin{tcolorbox}[colback=AIEPBlueSoft,colframe=AIEPCyan,boxrule=.6pt,arc=1mm]
\textbf{\color{AIEPNavy}Antes de avanzar:} una idea tecnológica comienza a tomar forma cuando el equipo puede nombrar
qué desea observar y para qué servirá ese dato.
\end{tcolorbox}

\newpage
% 2 ---------------------------------------------------------------------------
\pagehead{Del mundo físico al dato}{Definir variable, sensor y comportamiento del sistema}

\begin{fieldbox}[height=40mm]{Variable observable}
¿Qué cambia en el mundo real y podría medirse? Ejemplos: distancia, humedad, temperatura, luz, movimiento, sonido o presencia.
\writearea{13mm}
\end{fieldbox}

\begin{tabularx}{\textwidth}{@{}Y Y@{}}
\begin{fieldbox}[height=47mm]{Sensor elegido}
Nombre o tipo:\par\vspace{7mm}\hrule
\vspace{5mm}
¿Qué detecta?\par\vspace{7mm}\hrule
\end{fieldbox}
&
\begin{fieldbox}[height=47mm]{Por qué es coherente}
Expliquen la relación entre el sensor, la variable y el problema.
\writearea{16mm}
\end{fieldbox}
\end{tabularx}

\begin{accentfield}{La regla inicial}
{\large\bfseries\color{AIEPNavy}Cuando} \rule{63mm}{0.35pt}
\quad {\large\bfseries\color{AIEPRed}entonces} \rule{63mm}{0.35pt}
\vspace{5mm}
\end{accentfield}

\begin{fieldbox}[height=54mm]{El ciclo completo}
\begin{tabularx}{\textwidth}{@{}C{30mm}C{5mm}C{30mm}C{5mm}C{30mm}C{5mm}C{30mm}@{}}
\textbf{\color{AIEPNavy}ENTRADA} & $\rightarrow$ &
\textbf{\color{AIEPNavy}PROCESO} & $\rightarrow$ &
\textbf{\color{AIEPNavy}COMUNICACIÓN} & $\rightarrow$ &
\textbf{\color{AIEPNavy}ACCIÓN}\\
sensor y dato && regla o cálculo && LED, pantalla o app && alerta o respuesta\\
\end{tabularx}
\vspace{7mm}
\hrule
\vspace{8mm}
\hrule
\end{fieldbox}

\begin{fieldbox}[height=43mm]{Componentes que imaginamos}
Marquen y completen: \quad
\checkline{sensor}\quad
\checkline{placa}\quad
\checkline{salida}\quad
\checkline{software}\quad
\checkline{carcasa}\par
\vspace{6mm}
\textbf{¿Qué componente es indispensable para la primera prueba?}\par
\vspace{7mm}\hrule
\end{fieldbox}

\newpage
% 3 ---------------------------------------------------------------------------
\pagehead{Probar con criterio}{Diseñar una evidencia segura y verificable}

\begin{fieldbox}[height=48mm]{Primera pregunta de prueba}
¿Qué necesitan comprobar primero? Escríbanlo de modo que la respuesta pueda observarse.
\writearea{18mm}
\end{fieldbox}

\begin{tabularx}{\textwidth}{@{}Y Y@{}}
\begin{fieldbox}[height=62mm]{Procedimiento}
1. \rule{63mm}{0.35pt}\vspace{7mm}\par
2. \rule{63mm}{0.35pt}\vspace{7mm}\par
3. \rule{63mm}{0.35pt}\vspace{7mm}\par
4. \rule{63mm}{0.35pt}
\end{fieldbox}
&
\begin{fieldbox}[height=62mm]{Resultado esperado}
¿Qué dato o comportamiento mostraría que la idea funciona?
\writearea{25mm}
\end{fieldbox}
\end{tabularx}

\begin{accentfield}{Revisión de seguridad antes del USB}
\checkline{La placa está desconectada mientras conectamos.}\par
\checkline{Verificamos pinout, voltaje y polaridad.}\par
\checkline{Los LEDs tienen la resistencia correspondiente.}\par
\checkline{Existe una tierra común cuando corresponde.}\par
\checkline{Una persona responsable revisó el circuito antes de energizar.}
\end{accentfield}

\begin{fieldbox}[height=51mm]{Uso de IA y verificación}
\textbf{¿Qué le preguntamos?}\par\vspace{6mm}\hrule
\vspace{5mm}
\textbf{¿Qué dato, conexión o código verificamos?}\par\vspace{6mm}\hrule
\vspace{5mm}
\textbf{¿Qué aceptamos, corregimos o descartamos?}\par\vspace{6mm}\hrule
\end{fieldbox}

\begin{fieldbox}[height=36mm]{Evidencia obtenida o prueba pendiente}
\checkline{lectura}\quad
\checkline{simulación}\quad
\checkline{diagrama revisado}\quad
\checkline{código}\quad
\checkline{fotografía/video}\quad
\checkline{pendiente}\par
\writearea{8mm}
\end{fieldbox}

\newpage
% 4 ---------------------------------------------------------------------------
\pagehead{Convertir la idea en una ruta}{Roles, nivel actual, siguiente hito y relato breve}

\begin{fieldbox}[height=74mm]{Seis responsabilidades del equipo}
\begin{tabularx}{\textwidth}{@{}>{\bfseries\color{AIEPNavy}}p{52mm}Y@{}}
Coordinación & Nombre: \rule{91mm}{0.35pt}\\[3mm]
Problema y propósito & Nombre: \rule{91mm}{0.35pt}\\[3mm]
Investigación & Nombre: \rule{91mm}{0.35pt}\\[3mm]
Prototipo físico & Nombre: \rule{91mm}{0.35pt}\\[3mm]
Programación y datos & Nombre: \rule{91mm}{0.35pt}\\[3mm]
Relato y evidencia & Nombre: \rule{91mm}{0.35pt}\\
\end{tabularx}
\end{fieldbox}

\begin{fieldbox}[height=43mm]{Nivel de madurez actual}
\checkline{Idea definida}\quad
\checkline{Sensor elegido}\quad
\checkline{Circuito o simulación}\par
\checkline{Código}\quad
\checkline{Datos visibles}\quad
\checkline{Prototipo integrado}\par
\vspace{5mm}
\textbf{La evidencia que sostiene esta elección es:}\par\vspace{7mm}\hrule
\end{fieldbox}

\begin{accentfield}{Próximo hito verificable}
Antes de la siguiente instancia, el equipo tendrá:\par\vspace{4mm}\hrule
\vspace{5mm}
Se podrá comprobar mediante: \rule{92mm}{0.35pt}
\vspace{4mm}
\end{accentfield}

\begin{fieldbox}[height=54mm]{Pitch relámpago · 30 a 40 segundos}
\textbf{Problema:} observamos que \rule{120mm}{0.35pt}\par\vspace{5mm}
\textbf{Propuesta:} queremos usar \rule{106mm}{0.35pt}\par\vspace{5mm}
\textbf{Funcionamiento:} cuando \rule{46mm}{0.35pt}, entonces \rule{48mm}{0.35pt}\par\vspace{5mm}
\textbf{Siguiente paso:} ahora necesitamos comprobar \rule{87mm}{0.35pt}
\end{fieldbox}

\begin{tcolorbox}[colback=AIEPBlueSoft,colframe=AIEPCyan,boxrule=.6pt,arc=1mm]
\textbf{\color{AIEPNavy}Salida del taller:}\quad
\checkline{problema}\quad
\checkline{variable}\quad
\checkline{sensor}\quad
\checkline{regla}\quad
\checkline{prueba segura}\quad
\checkline{siguiente hito}
\end{tcolorbox}

\begin{tcolorbox}[colback=AIEPNavy,colframe=AIEPNavy,arc=1mm,halign=center]
{\large\bfseries\color{white}Una buena idea no se adivina: se observa, se prueba y se mejora.}
\end{tcolorbox}

\end{document}
